%%#! latex

%%%%%%%%%%             Preamble                   %%%%%%%%%%

\documentclass[10pt]{jarticle}
\usepackage{amsfonts,amssymb,amsmath,amscd,amsthm,latexsym}
\usepackage{graphicx}
\usepackage{fancybox}
\usepackage{okumacro}
%\usepackage{hyperref}
%\usepackage{showkeys}
%\usepackage{fancybox,ascmac}
%\usepackage{fancyhdr}
\pagestyle{empty}

%%%%%%%%%%             Layout                     %%%%%%%%%%

\setlength{\topmargin}{-1cm}
\setlength{\textheight}{25cm}
\setlength{\textwidth}{16cm}
%\setlength{\oddsidemargin}{0cm}
%\setlength{\evensidemargin}{0cm}
%\addtolength{\textwidth}{.815in}
\addtolength{\oddsidemargin}{-2cm}
\addtolength{\evensidemargin}{-2cm}

%%%%%%%%%%           Definition, Theorem, etc.     %%%%%%%%%%

\newtheorem{thm}{定理}
\newtheorem{lem}[thm]{補題}
\newtheorem{prop}[thm]{命題}
\newtheorem{claim}[thm]{主張}
\newtheorem{cor}[thm]{系}
\newtheorem{prob}[thm]{問題}
\newtheorem{conj}[thm]{予想}
\theoremstyle{definition}
\newtheorem{dfn}[thm]{定義} 
\newtheorem{ex}[thm]{例}
\newtheorem{reidai}[thm]{例題}
\newtheorem{rmk}[thm]{注意}

%\numberwithin{thm}{section}
%\numberwithin{equation}{section}

%%%%%%%%%%             Macro (by Nasu)           %%%%%%%%%%

\newcommand{\Proof}{\noindent {\bf 証明)}\ \ }
\newcommand{\Hom}{\operatorname{Hom}}
\newcommand{\Ext}{\operatorname{Ext}}
\newcommand{\Spec}{\operatorname{Spec}}
\newcommand{\Proj}{\operatorname{Proj}}
\newcommand{\ob}{\operatorname{ob}}
\newcommand{\Hilb}{\operatorname{Hilb}}
\newcommand{\red}{\operatorname{red}}
\newcommand{\im}{\operatorname{im}}
\newcommand{\Bs}{\operatorname{Bs}}
\newcommand{\Jac}{\operatorname{Jac}}
\newcommand{\Pic}{\operatorname{Pic}}
\newcommand{\coker}{\operatorname{coker}}
\newcommand{\Gr}{\operatorname{Gr}}
\newcommand{\Bl}{\operatorname{Bl}}
%\newcommand{\algeeq}{\stackrel{alg.}{\; \sim \;}}
\newcommand{\Rat}{\operatorname{Rat}}
\newcommand{\Exc}{\operatorname{Exc}}
\newcommand{\car}{\operatorname{char}}
\newcommand{\Aut}{\operatorname{Aut}}
\newcommand{\End}{\operatorname{End}}
\newcommand{\sat}{\operatorname{sat}}
\newcommand{\pd}{\operatorname{dp}}
\newcommand{\rank}{\operatorname{rank}}
\newcommand{\Ann}{\operatorname{Ann}}

\renewcommand{\theenumi}{\arabic{enumi}}
\renewcommand{\theenumii}{\alph{enumii}}
\renewcommand{\theenumiii}{\roman{enumiii}}
\renewcommand{\labelenumi}{(\theenumi)}
\renewcommand{\labelenumii}{(\theenumii)}
\renewcommand{\labelenumiii}{[\theenumiii]}
%\renewcommand{\baselinestretch}{1.2}
%\renewcommand{\abstractname}{\sc Abstract}

%%%%%%%%%%           Commutative diagram          %%%%%%%%%%

% A litte smaller diagram than AMS CD envrionment

\newcommand{\mapright}[1]{%
\smash{\mathop{%
\hbox to 1cm{\rightarrowfill}}\limits^{#1}}}
\newcommand{\mapleft}[1]{%
\smash{\mathop{%
\hbox to 1cm{\leftarrowfill}}\limits^{#1}}}
\newcommand{\mapdown}[1]{\Big\downarrow
\llap{$\vcenter{\hbox{$\scriptstyle#1\,$}}$ }}
\newcommand{\mapup}[1]{\Big\uparrow
\rlap{$\vcenter{\hbox{$\scriptstyle#1$}}$ }}

%%%%%%%%%%              Title                    %%%%%%%%%%

\title{}
\author{}
\date{}

%%%%%%%%%%            Text Start                 %%%%%%%%%%

\begin{document}

%\maketitle

\begin{center}
 {\bf \Huge 東海大学理学部数学・情報数理談話会}
\end{center}

\bigskip

{\Large
以下の要領において談話会を開催致します. 多数の方の御来聴を
お待ち致しております. 

\bigskip

\noindent
日時: 2026年5月12日(火) 16:00--17:00\\
\smallskip
会場: 東海大学湘南校舎18号館8階理学部ゼミ室3 (18-831)\\

\medskip

\noindent
講演者: 
Euisung Park (Department of Mathematics, Korea University)\\
\smallskip
タイトル: 
Secant Rank and Syzygies of Projections of Elliptic Normal Curves\\

\noindent
アブストラクト:
In this talk, we study the syzygies of projections of elliptic normal curves. This talk is based on joint work with Changho Han. Let
$C \subset \mathbb{P}^{d-1}$ be an elliptic normal curve and let $C_q$ denote the projection of $C$ from a point $q$. A natural question is how the syzygetic properties of the projected curve depend on the position of the projection center.
Using the geometry of secant varieties, we relate the Green--Lazarsfeld index of $C_q$ to the secant rank of $q$ with respect to $C$. In particular, we show that if $q$ lies in the $s$-th secant variety of $C$, then $\mathrm{index}(C_q) \leq s-3$, and equality holds for a general point in that secant variety.
Our approach realizes projected elliptic curves as hyperplane sections of suitable elliptic ruled surface scrolls and studies their minimal free resolutions.

\vskip 2cm

\begin{flushright}
  世話人: 那須 弘和\\
  nasu@tokai.ac.jp
\end{flushright}
}
\end{document}

